% All rights to the corporate identity of University of Amsterdam and its elements
% belong to the University of Amsterdam.
% Please refer to the following page for further information:
% https://extranet.uva.nl/en/a-z/content/corporate-communication/corporate-communication.html
%
% Copyright (c) 2015 Franz Geiger
%
% Permission is hereby granted, free of charge, to any person obtaining a copy
% of this software and associated documentation files (the "Software"), to deal
% in the Software without restriction, including without limitation the rights
% to use, copy, modify, merge, publish, distribute, sublicense, and/or sell
% copies of the Software, and to permit persons to whom the Software is
% furnished to do so, subject to the following conditions:
%
% The above copyright notice and this permission notice shall be included in
% all copies or substantial portions of the Software.
%
% THE SOFTWARE IS PROVIDED "AS IS", WITHOUT WARRANTY OF ANY KIND, EXPRESS OR
% IMPLIED, INCLUDING BUT NOT LIMITED TO THE WARRANTIES OF MERCHANTABILITY,
% FITNESS FOR A PARTICULAR PURPOSE AND NONINFRINGEMENT.  IN NO EVENT SHALL THE
% AUTHORS OR COPYRIGHT HOLDERS BE LIABLE FOR ANY CLAIM, DAMAGES OR OTHER
% LIABILITY, WHETHER IN AN ACTION OF CONTRACT, TORT OR OTHERWISE, ARISING FROM,
% OUT OF OR IN CONNECTION WITH THE SOFTWARE OR THE USE OR OTHER DEALINGS IN
% THE SOFTWARE.
%
% Additionally, this example document may be used with the template without
% above license notice.

\documentclass[pdf]{beamer}
% Load University of Amsterdam theme
% Options:
% 	logo: Select whether to use the Dutch or English logo in header.
% 		Either 'NL' or 'ENG'. Default: ENG
% 	maincolor: Select which main color to use for title page and as highlight color.
%		Possible values are: Any color definition. Default: uva-rood
% 		This theme pre-defines colors from the University of Amsterdam corporate identity:
% 			uva-zwart	(general UvA Black)
% 			uva-rood	(general UvA Red)
% 			uva-feb-groen	(Green for Faculty of Economics and Business)
% 			uva-fgw-blauw	(Blue for Faculty of Humanities)
% 			uva-fdr-donker-rood	(Dark Red for Faculty of Law)
% 			uva-gen-groen	(Green for Faculty of Medicine (AMC))
% 			uva-fnwi-paars	(Purple for Faculty of Science)
% 			uva-fmg-oranje	(Orange for Faculty of Social and Behavioural Sciences)
% 			uva-iis-licht-blauw	(Light Blue for Institute of Interdisciplinary Studies)
% 			uva-ilo-rood	(Red for Interfacultaire Lerarenopleidingen)
% 		Please refer to the following page for usage instructions regarding the
% 		corporate identity colors: https://extranet.uva.nl/en/a-z/content/corporate-communication/corporate-identity-elements/colour/colour-kopie.html
\usetheme[logo=ENG, maincolor=uva-rood]{UniversityOfAmsterdam}

\usepackage[utf8]{inputenc}
\usepackage{hyperref}

\author{Author}
\title{Title}
\subtitle{Subtitle}
\institute{Institute}
\date{\today}

\begin{document}

% Title page (plain in order to hide regular background and footline)
\frame[plain]{\titlepage}

\begin{frame}
	\frametitle{Outline}
	\tableofcontents
\end{frame}

\section{Select your faculty color scheme}

\begin{frame}{Pre-defined colors from UvA corporate identity}
	\newcounter{posX}
	\begin{tikzpicture}
		\edef\posX{0}
		\edef\posY{11}
		\foreach \color in {uva-zwart, uva-rood, uva-feb-groen, uva-fgw-blauw, uva-fdr-donker-rood, uva-gen-groen, uva-fnwi-paars, uva-fmg-oranje, uva-iis-licht-blauw, uva-ilo-rood} {
			\node [fill=\color, text=white] (\color) at (\posX, \posY) {\color};

			\pgfmathparse{\posX+1}
			\xdef\posX{\pgfmathresult}
			\ifnum8>\posY
				\xdef\posY{11}
			\else
				\pgfmathparse{\posY-1}
				\xdef\posY{\pgfmathresult}			
			\fi
			};
	\end{tikzpicture}
\end{frame}

\section{Example Content}

\begin{frame}{\insertsectionhead}{Some lists}
	\begin{columns}
		\begin{column}{0.5\textwidth}
			\begin{enumerate}
				\item Foo
				\item Bar
				\begin{enumerate}
					\item test
					\item hello world
					\begin{enumerate}
						\item more levels
					\end{enumerate}
				\end{enumerate}
				\item Baz
			\end{enumerate}
		\end{column}
		\begin{column}{0.5\textwidth}
			\begin{itemize}
				\item Foo
				\item Bar
				\begin{itemize}
					\item test
					\item hello world
					\begin{itemize}
						\item more levels
					\end{itemize}
				\end{itemize}
				\item Baz
			\end{itemize}
		\end{column}
	\end{columns}
\end{frame}

{ % Use a frametitle with title and subtitle on one line to safe space
	\setbeamertemplate{frametitle}[University of Amsterdam oneline]

	\begin{frame}{Block elements}{With subtitle on the same line as title}
		\begin{columns}[t]
			\begin{column}{0.45\textwidth}
				\begin{example}[The usual]
					Hello World
				\end{example}
				\begin{theorem}[Stuff is hard]
					content...
				\end{theorem}
				\begin{corollary}
					content...
				\end{corollary}
			\end{column}
			\begin{column}{0.45\textwidth}
				\begin{fact}
					content...
				\end{fact}
				\begin{lemma}
					content...
				\end{lemma}
				\begin{proof}[Proof \alert{this!}]
					trivial
				\end{proof}
			\end{column}
		\end{columns}
	\end{frame}
}

\begin{frame}{Even more boxes}

	\begin{definition}[of a definition]
		A \alert{definition} is what is in such a box.
	\end{definition}
	\begin{block}{Title of the block element}
		Just another block with content.
	\end{block}
	\begin{alertblock}{Attention!}
		content...
	\end{alertblock}
	
\end{frame}

\section{Plain frames are really boring}

\begin{frame}[plain]

	A plain frame with a figure and a table.

	\vfill

	\begin{figure}
		\caption{University of Amsterdam}\label{fig:uva-logo}
		\includegraphics{beamerthemeUniversityOfAmsterdam/logo/uva_logo_ENG.pdf}
	\end{figure}
	
	\begin{table}
		\begin{tabular}{|c|c|c|}
			\hline
			A & B & C \\
			\hline
			D & E & F \\
			\hline
		\end{tabular}
		\caption{A table without content}\label{tab:nonsense}
	\end{table}
\end{frame}

\section{Section titles show up in the footer}

\begin{frame}{\insertsectionhead}
	\begin{alertblock}{And what if there is no section?}
		If a frame is outside any sections, (short) title of the presentation if used. In case this is empty too, author is used.
	\end{alertblock}
\end{frame}
\end{document}